\documentclass[11pt]{article}
\usepackage[margin=1in]{geometry}
\usepackage{booktabs}
\usepackage{hyperref}
\usepackage{xcolor}
\usepackage{natbib}
% Generated by src/make_macros.py from results/stats.json. Do not edit.
\newcommand{\PoneBuildsHundred}{7}
\newcommand{\PoneBuildsBlackwell}{2}
\newcommand{\PoneRowsPerMode}{38{,}496}
\newcommand{\PoneBadTotal}{0}
\newcommand{\PoneComparisons}{769{,}920}
\newcommand{\PoneCtlMinH}{6{,}784}
\newcommand{\PoneCtlMaxH}{8{,}544}
\newcommand{\PoneCtlB}{19{,}544}
\newcommand{\PoneCublasVersions}{4}
\newcommand{\PtwoConfigs}{7}
\newcommand{\PtwoPassConfigs}{6}
\newcommand{\PtwoCtlDivMin}{4\%}
\newcommand{\PtwoCtlDivMax}{76\%}
\newcommand{\PtwoCostCells}{17}
\newcommand{\PtwoCostPass}{16}
\newcommand{\PtwoPrefillPlainMax}{5.0\%}
\newcommand{\PtwoPrefillStockMin}{9.5\%}
\newcommand{\PtwoPrefillStockMax}{25.5\%}
\newcommand{\QsevenStock}{63.1\%}
\newcommand{\QsevenPlain}{0.2\%}
\newcommand{\QsevenPlainMax}{3.1\%}
\newcommand{\GemmaStock}{59.2\%}
\newcommand{\GemmaPlain}{1.0\%}
\newcommand{\TwoPStock}{61.0\%}
\newcommand{\TwoPPlain}{5.8\%}
\newcommand{\BigStock}{58.2\%}
\newcommand{\BigPlain}{$-$2.3\%}
\newcommand{\TwoPComm}{6.0\%}
\newcommand{\DiagReq}{30}
\newcommand{\DiagTritonOn}{0}
\newcommand{\DiagTritonOff}{19}
\newcommand{\DiagEagerOn}{0}
\newcommand{\DiagEagerOff}{23}
\newcommand{\DiagNoscOn}{15}
\newcommand{\DiagNoscOff}{24}
\newcommand{\DiagSevenOn}{0}
\newcommand{\DiagSevenOff}{30}
\newcommand{\DiagNoswOn}{0}
\newcommand{\DiagNoswOff}{17}
\newcommand{\DiagNocgOn}{0}
\newcommand{\DiagNocgOff}{15}
\newcommand{\ShapeCount}{37}
\newcommand{\ShapeMaxFailN}{96}
\newcommand{\ShapeMinPassN}{128}
\newcommand{\RouterBfFirstBad}{23}
\newcommand{\RouterFpFirstBad}{2}
\newcommand{\MidLo}{97}
\newcommand{\MidHi}{127}
\newcommand{\MidMult}{8}
\newcommand{\NarrowNs}{64 and 96}
\newcommand{\NarrowFailK}{2880, 4096 and 5120}
\newcommand{\NarrowPassK}{512, 896, 1024, 1536, 2048 and 8192}
\newcommand{\NoenvFailK}{2880, 4096 and 5120}
\newcommand{\NoenvMinFailN}{64}
\newcommand{\NoenvMaxFailN}{1024}
\newcommand{\NoenvMinPassN}{2048}
\newcommand{\MistralStock}{5.8\%}
\newcommand{\MistralPlain}{3.3\%}
\newcommand{\GemmaOnDiv}{40}
\newcommand{\GemmaPlainDiv}{39}
\newcommand{\GemmaOffDiv}{60}
\newcommand{\GemmaReq}{100}
\newcommand{\OffDivQsevenMain}{41}
\newcommand{\OffReqQsevenMain}{100}
\newcommand{\StockMainQseven}{62.3\%}
\newcommand{\PlainMainQseven}{0.7\%}
\newcommand{\TorchDefaultBad}{1{,}200}
\newcommand{\TorchDefaultRows}{12{,}200}
\newcommand{\LlamaStock}{$-$0.5\%}
\newcommand{\QthreeStock}{9.1\%}
\newcommand{\HtwoShapes}{96}
\newcommand{\HtwoAgree}{96}
\newcommand{\HtwoAgreeNoenv}{96}
\newcommand{\GptossTests}{4}
\newcommand{\GptossOnPass}{4}
\newcommand{\GptossVtwoPass}{4}
\newcommand{\GptossBlanketPass}{0}
\newcommand{\FafixBase}{5}
\newcommand{\FafixFix}{0}
\newcommand{\FafixN}{20}
\newcommand{\TuneShapes}{5}
\newcommand{\TuneRows}{146{,}010}
\newcommand{\TuneBad}{0}
\newcommand{\TuneCrossReq}{30}
\newcommand{\TuneCrossMis}{0}
\newcommand{\TuneStockB}{64.4\%}
\newcommand{\TunedB}{10.0\%}
\newcommand{\TuneLossMin}{5.1\%}
\newcommand{\TuneLossMax}{11.7\%}
\newcommand{\TuneStockMin}{33.3\%}
\newcommand{\TuneStockMax}{67.7\%}
\newcommand{\TuneBmax}{128}
\newcommand{\TuneSpeedMin}{2.10}
\newcommand{\TuneSpeedMax}{3.37}
\newcommand{\TuneLargeMin}{1.01}
\newcommand{\TuneLargeMax}{1.51}
\newcommand{\TuneSmallMin}{0.97}
\newcommand{\TuneSmallMax}{1.15}
\newcommand{\OneEVerdict}{\pending{p1e.verdict}}
\newcommand{\OneEShapes}{\pending{p1e.C.shapes}}
\newcommand{\OneEBadC}{\pending{p1e.C.bad}}
\newcommand{\OneECompC}{\pending{p1e.C.comparisons}}
\newcommand{\OneEErrC}{\pending{p1e.C.error-shapes}}
\newcommand{\OneEBadB}{\pending{p1e.B.bad}}
\newcommand{\OneEBadShapesB}{\pending{p1e.B.bad-shapes}}
\newcommand{\OneEBadD}{\pending{p1e.D.bad}}
\newcommand{\OneEBadE}{\pending{p1e.E.bad}}
\newcommand{\OneEErrA}{\pending{p1e.A.error-shapes}}
\newcommand{\ServeUnit}{40}
\newcommand{\ServeRounds}{\pending{serve.decode.rounds}}
\newcommand{\ServeWins}{\pending{serve.decode.tuned-wins}}
\newcommand{\ServeDecGain}{\pending{serve.decode.gain-over-stock}}
\newcommand{\ServeDecStockLoss}{\pending{serve.decode.stock-loss-vs-off}}
\newcommand{\ServeDecTunedLoss}{\pending{serve.decode.tuned-loss-vs-off}}
\newcommand{\ServePreGain}{\pending{serve.prefill.gain-over-stock}}
\newcommand{\ServePreTunedLoss}{\pending{serve.prefill.tuned-loss-vs-off}}

\newcommand{\todo}[1]{\textcolor{red}{[#1]}}

\title{Batch-Invariant LLM Serving on Hopper Without Custom Kernels}
\author{Isaac M. Thompson\\Brigham Young University}
\date{}

\begin{document}
\maketitle

\begin{abstract}
Greedy decoding in a batched inference server is not deterministic: the same request can produce different tokens depending on which other requests share its batch and whether its prefix is cached. vLLM's batch-invariant mode removes this dependence, but on an H100 it cuts decode throughput at batch 32 by \QsevenStock{} for a model without hand-tuned kernels. We show that on Hopper and Blackwell GPUs much of this cost can be avoided. vLLM routes Linear layers through a slower Triton kernel in this mode, a choice made after a model failed batch invariance with cuBLAS. Under the mode's settings, however, cuBLAS is batch-invariant for the large Linear shapes we tested: across \PoneBuildsHundred{} H100 builds, \PoneBuildsBlackwell{} B200 builds and \PoneCublasVersions{} cuBLAS versions, \PoneBadTotal{} of \PoneComparisons{} row comparisons differ under batch-invariant mode's flags. Routing Linear layers to cuBLAS gives identical tokens and logprobs in \PtwoPassConfigs{} of \PtwoConfigs{} serving configurations, including tensor parallelism, long context with chunked prefill and seeded sampling, while cutting the throughput loss to \QsevenPlain{} for Qwen2.5-7B and \GemmaPlain{} for Gemma-2-9B. The remaining configuration, Gemma-2-9B, is not batch-invariant under vLLM's own mode either; we localize that failure to FlashAttention's sliding-window path under CUDA graphs. At two-way tensor parallelism, the residual overhead comes from batch-invariant communication settings rather than the routing. Narrow layers such as mixture-of-experts routers can still fail with cuBLAS, so the routing has to be decided per shape.
\end{abstract}

\section{Introduction}
Greedy decoding is often assumed to be deterministic, but in a batched inference server it is not. The same request, served alone or alongside unrelated requests, with or without a cached prefix, can produce different tokens, because the floating-point reduction order of matrix multiplication and attention kernels depends on the shape of the batch \citep{he2025nondeterminism,yuan2025numerical}. In our measurements on vLLM \citep{kwon2023vllm}, \OffDivQsevenMain{} of \OffReqQsevenMain{} greedy math requests to Qwen2.5-7B-Instruct changed tokens across these conditions.

vLLM offers a batch-invariant mode that removes this dependence by replacing reduction kernels with fixed-order implementations. On an H100 it reduced decode throughput at batch 32 by \StockMainQseven{} for Qwen2.5-7B-Instruct. vLLM has since added hand-tuned kernel configurations for a few models \citep{vllm62219}, which recover most of the loss for those models only.

On Hopper (SM90) and Blackwell (SM100) GPUs, vLLM's batch-invariant mode sends every unquantized Linear layer through its Triton kernel, a choice made after GPT-OSS failed batch invariance with cuBLAS on H100 \citep{vllm35404}. We find that for the large Linear shapes that dominate inference cost, cuBLAS is batch-invariant under the mode's settings and even produces the same bits as the Triton kernel. Routing these layers to cuBLAS keeps outputs bitwise identical across serving conditions and removes nearly all of the cost. Narrow shapes, such as a mixture-of-experts router, are not covered by this finding. We contribute:
\begin{itemize}
\item a direct bitwise test showing that cuBLAS matrix multiplication is batch-invariant for the tested shapes under batch-invariant mode's flags on H100 and B200, across \PoneCublasVersions{} cuBLAS versions (\PoneComparisons{} row comparisons, \PoneBadTotal{} mismatches);
\item end-to-end serving tests over tensor parallelism, long context with chunked prefill, seeded sampling, and four model families, with preregistered pass/fail rules and controls;
\item throughput measurements against vLLM's own batch-invariant mode on current vLLM main;
\item a diagnosis of the two places it falls short: the TP=2 overhead comes from fixed-order communication, and a Gemma-2 failure lies in vLLM's FlashAttention path, not in the routing.
\end{itemize}

\section{Background}
A batch-invariant kernel computes each row of its output with the same sequence of floating-point operations regardless of how many other rows are in the batch. Reductions whose split depends on batch size (split-K matrix multiplication, split-KV attention, and some normalization kernels) break this property \citep{he2025nondeterminism}. vLLM's batch-invariant mode, enabled with \texttt{VLLM\_BATCH\_INVARIANT=1}, pins attention to a single split, replaces reductions with fixed-order kernels, disables reduced-precision reductions, prefers cuBLASLt, and pins collective communication to fixed algorithms. On Ampere (SM80) it overrides matrix multiplication with a persistent Triton kernel, since cuBLASLt there cannot be made batch-invariant. On SM90 and SM100 it instead sets cuBLAS workspace variables, with a source comment stating that split-K is the only source of batch variance on those GPUs. That comment is not accurate for every shape: after GPT-OSS failed batch invariance with cuBLAS on H100, \texttt{UnquantizedLinearMethod.apply} was changed to call the Triton \texttt{linear\_batch\_invariant} on every CUDA device \citep{vllm35404}. On Hopper the Linear layers, which perform most of the matrix-multiplication work, therefore run the Triton kernel.

\section{Method}
The change is one routing decision: under batch-invariant mode on SM90 and newer, Linear layers call \texttt{torch.nn.functional.linear} (cuBLAS) instead of the Triton kernel. Everything else in batch-invariant mode is unchanged. The routing applies to bf16 and fp16 inputs only: in the same bitwise test, cuBLAS fp32 matrix multiplication was not batch-invariant (rows differed at every batch size above one), while bf16 and fp16 with a bias term were. Our prototype applies this as an import hook that patches every vLLM process, including tensor-parallel workers, and logs a marker per process; a run counts only if the expected markers are present. An upstream version should also run a short bitwise self-test at startup and fall back to the Triton path if cuBLAS ever fails it, since NVIDIA does not document this behavior as a guarantee.

\section{Evaluation}
\paragraph{Setup.} vLLM main (0.33.0rc1.dev454), PyTorch 2.14 with CUDA 13.0, bf16, NVIDIA H100 80GB unless noted. Determinism: each request is served 8 times in random order under four conditions (cold or warm prefix cache, alone or with 7 unrelated requests, each twice), and every pair of the 8 outputs is compared token by token; top-20 logprobs are also recorded. Prompts are few-shot GSM8K and MATH-500 problems. Pass/fail rules were written before each run (repository \texttt{PROTOCOL.md}); a control with batch-invariant mode off must diverge for a test to count.

\paragraph{Is cuBLAS batch-invariant?} We multiply each row alone and inside batches of 1 to 8192 rows, at every Linear shape of six models, in bf16 and fp16, and compare rows near the start and end of the batch (Table~\ref{tab:bitwise}). Under batch-invariant mode's flags, no row differed on any of \PoneBuildsHundred{} H100 builds or \PoneBuildsBlackwell{} B200 builds, including with a second process saturating the GPU. Setting only the PyTorch flags, without the cuBLAS environment variables, made no difference. With PyTorch defaults, between \PoneCtlMinH{} and \PoneCtlMaxH{} of \PoneRowsPerMode{} rows differed on H100, and \PoneCtlB{} on B200, so the test detects batch dependence when it exists.

\begin{table}[t]\centering\small
% Generated by src/make_macros.py. Do not edit.
\begin{tabular}{lrrr}
\toprule
GPU, torch build & BI flags & Flags only & Defaults \\
\midrule
B200, torch t2.13.0+cu130 & 0 & 0 & 19{,}544 \\
B200, torch t2.14.0+cu130 & 0 & 0 & 19{,}544 \\
H100, torch t2.11.0+cu128 & 0 & 0 & 7{,}344 \\
H100, torch t2.12.0+cu126 & 0 & 0 & 6{,}784 \\
H100, torch t2.13.0+cu126 & 0 & 0 & 6{,}784 \\
H100, torch t2.13.0+cu129 & 0 & 0 & 7{,}344 \\
H100, torch t2.13.0+cu130 & 0 & 0 & 8{,}544 \\
H100, torch t2.14.0+cu126 & 0 & 0 & 6{,}784 \\
H100, torch t2.14.0+cu130 & 0 & 0 & 8{,}544 \\
H100, torch t2.14.0+cu130 + GPU load & 0 & 0 & 8{,}544 \\
\bottomrule
\end{tabular}

\caption{Rows that differ between a row multiplied alone and inside a batch, out of \PoneRowsPerMode{} per cell. BI flags: environment variables and PyTorch flags as batch-invariant mode sets them. Flags only: PyTorch flags without the cuBLAS environment variables. Defaults: PyTorch defaults (control).}\label{tab:bitwise}
\end{table}

\paragraph{End-to-end determinism.} With Linear layers on cuBLAS, \PtwoPassConfigs{} of \PtwoConfigs{} serving configurations produced identical tokens and identical top-k logprobs across all conditions (Table~\ref{tab:det}). Controls with batch-invariant mode off diverged on \PtwoCtlDivMin{} to \PtwoCtlDivMax{} of requests. Gemma-2-9B fails: \GemmaPlainDiv{} of \GemmaReq{} requests diverged with cuBLAS routing, and \GemmaOnDiv{} with vLLM's stock batch-invariant mode, against \GemmaOffDiv{} with the mode off. vLLM's batch-invariant mode does not cover Gemma-2 in this setup; the routing change neither causes nor fixes this. Varying one component at a time under stock batch-invariant mode (\DiagReq{} requests each) localizes the failure: switching attention to the Triton backend removes it (\DiagTritonOn{} diverging requests, against \DiagTritonOff{} for its control), as does eager execution (\DiagEagerOn{}, against \DiagEagerOff{}). Removing logit softcapping does not (\DiagNoscOn{}, against \DiagNoscOff{}), and Gemma-7B, which shares Gemma-2's head size but has no softcapping, is unaffected (\DiagSevenOn{}, against \DiagSevenOff{}). Two further variants narrow it: disabling Gemma-2's sliding-window attention removes the divergence (\DiagNoswOn{}, against \DiagNoswOff{}), and so does turning off CUDA graphs while keeping compilation (\DiagNocgOn{}, against \DiagNocgOff{}). The failure therefore lies in FlashAttention's sliding-window path under CUDA graphs.

\begin{table}[t]\centering\small
% Generated by src/make_macros.py. Do not edit.
\begin{tabular}{lrrr}
\toprule
Configuration & Plain cuBLAS & Mode off & Markers \\
\midrule
Qwen2.5-7B, TP=2 & 0 / 100 & 38 / 100 & 2 \\
Qwen2.5-7B, TP=4 (TP=2 control) & 0 / 100 & 38 / 100 & 4 \\
Qwen2.5-7B, long context, chunked prefill & 0 / 50 & 29 / 50 & 1 \\
Qwen2.5-7B, seeded sampling & 0 / 100 & 40 / 100 & 1 \\
Qwen2.5-32B, TP=2 & 0 / 100 & 76 / 100 & 2 \\
Mistral-7B & 0 / 100 & 4 / 100 & 1 \\
Gemma-2-9B & 39 / 100 & 60 / 100 & 1 \\
\bottomrule
\end{tabular}

\caption{Requests whose 8 runs were not all identical, with Linear layers on cuBLAS under batch-invariant mode (plain cuBLAS) and with batch-invariant mode off (control). Markers count processes that loaded the routing change.}\label{tab:det}
\end{table}

\paragraph{Throughput.} Table~\ref{tab:cost} reports decode throughput loss at batch 32 against batch-invariant mode off, the median of three independent jobs. For models without tuned configurations, stock batch-invariant mode loses \QsevenStock{} (Qwen2.5-7B) and \GemmaStock{} (Gemma-2-9B); cuBLAS routing loses \QsevenPlain{} and \GemmaPlain{}. Mistral-7B shares most Linear shapes with Llama-3.1-8B, which vLLM has tuned, so stock mode is already cheap there (\MistralStock{}). For prefill, cuBLAS routing loses at most \PtwoPrefillPlainMax{} against \PtwoPrefillStockMin{} to \PtwoPrefillStockMax{} for stock mode. \PtwoCostPass{} of \PtwoCostCells{} measured cells meet our preregistered limit of 5\%. The exception is Qwen2.5-7B at TP=2, which loses \TwoPPlain{} (stock: \TwoPStock{}). This excess comes from batch-invariant mode's fixed-order communication settings, not from the routing change: applying only those settings, with batch-invariant mode otherwise off, costs \TwoPComm{} on its own. Cheaper deterministic communication is a separate target.

\begin{table}[t]\centering\small
% Generated by src/make_macros.py. Do not edit.
\begin{tabular}{lrrr}
\toprule
Model & Stock & Plain cuBLAS & Plain, worst repeat \\
\midrule
Qwen2.5-7B & 63.1\% & 0.2\% & 3.1\% \\
Gemma-2-9B & 59.2\% & 1.0\% & 1.3\% \\
Mistral-7B & 5.8\% & 3.3\% & 5.9\% \\
Qwen2.5-7B, TP=2 & 61.0\% & 5.8\% & 6.8\% \\
Qwen2.5-32B, TP=2 & 58.2\% & $-$2.3\% & $-$1.0\% \\
\bottomrule
\end{tabular}

\caption{Decode throughput loss at batch 32 against batch-invariant mode off. Median of three jobs.}\label{tab:cost}
\end{table}


\section{Limitations}
The bitwise property is observed, not guaranteed by NVIDIA, and it does not hold for every shape: another contributor found a mixture-of-experts router shape (N=32, K=2880) that is not batch-invariant under the same settings, so an upstream change must decide per shape, for example with a bitwise check of each Linear shape at model load. We tested one engine, single-node tensor parallelism up to four GPUs, dense models up to 32B parameters, and bf16; we did not run serving tests on H200, multi-node serving, mixture-of-experts models, quantized models, or production traffic. Gemma-2-9B is not batch-invariant under vLLM's batch-invariant mode with its default FlashAttention backend. For models in vLLM's tuned table the benefit is small. We compare only against vLLM's own batch-invariant mode, not against other engines' deterministic modes or LLM-42 \citep{gond2026llm42}.

\section{Related Work}
\citet{he2025nondeterminism} identify batch-size-dependent reductions as the main source of inference nondeterminism and describe batch-invariant kernels; vLLM's batch-invariant mode follows this design. \citet{yuan2025numerical} show that numerical nondeterminism changes reasoning accuracy across batch sizes, GPUs, and precisions. LLM-42 \citep{gond2026llm42} obtains determinism through verified speculation rather than invariant kernels. Our contribution is narrower: on current NVIDIA hardware, under the right flags, the vendor matrix-multiplication library is batch-invariant for the large Linear shapes that dominate cost, so the custom kernel is unnecessary for those shapes. It remains necessary for others, such as small mixture-of-experts routers.

\section{Conclusion}
Reproducible serving has been treated as a trade between determinism and throughput. On current NVIDIA hardware, that trade is mostly unnecessary for dense Linear layers: the vendor library already provides batch invariance under the right flags, and routing to it recovers nearly all of the lost throughput. We hope this makes batch-invariant serving a reasonable default for evaluation and for reinforcement-learning rollouts, where reproducibility matters most. An upstream version of the change should keep a startup self-test, since the property is observed rather than guaranteed. Code, raw outputs, and the decision rules for every experiment are at \url{https://github.com/imthomp/vllm-batch-invariance}.

\bibliographystyle{plainnat}
\bibliography{refs}
\end{document}
